\documentclass[10pt,a4paper]{amsart}
\usepackage[margin=19mm]{geometry}
\usepackage{amsmath,amssymb,mathtools}
\usepackage[scr=rsfs,cal=euler]{mathalfa}
\usepackage{xfrac}
\usepackage[unicode,hidelinks,bookmarks=false]{hyperref}
\newcommand{\ie}{i.e.\ }
\newcommand{\cf}{cf.\ }
\newcommand{\resp}{respectively\ }
\newcommand{\Subsection}{\subsection}
\providecommand{\nobreakdash}{\nobreak\hbox{-}}
\setlength{\emergencystretch}{2em}
\multlinegap0pt
\usepackage{amssymb}
\newtheorem{lem}{Lemma}
\newtheorem{prop}{Proposition}
\newtheorem{cor}{Corollary}
\newcommand{\Ac}{\mathsf{Ac}}
\newcommand{\CA}{\mathsf{CA}}
\newcommand{\Com}{\mathsf{Com}}
\DeclareMathOperator{\Ext}{Ext}
\newcommand{\Flat}{\mathsf{Flat}}
\DeclareMathOperator{\Hom}{Hom}
\newcommand{\Hot}{\mathsf{Hot}}
\newcommand{\VF}{\mathsf{VF}}
\newcommand{\bctr}{\mathsf{bctr}}
\newcommand{\boZ}{\mathbb Z}
\newcommand{\bu}{{\text{\smaller\smaller$\scriptstyle\bullet$}}}
\newcommand{\inj}{\mathsf{inj}}
\newcommand{\proj}{\mathsf{proj}}
\newcommand{\rarrow}{\longrightarrow}
\newcommand{\sA}{\mathsf A}
\newcommand{\sB}{\mathsf B}
\newcommand{\sE}{\mathsf E}
\newcommand{\sK}{\mathsf K}
\newcommand{\sS}{\mathsf S}
\title{Coderived and contraderived categories \\ for a cotorsion pair, flat-type cotorsion pairs, \\ and relative periodicity}
\author{Leonid Positselski}
\date{Source: arXiv:2509.07645v4 (CC BY 4.0)}
\begin{document}
\maketitle

\noindent\emph{Source and licence.} Coderived and contraderived categories \\ for a cotorsion pair, flat-type cotorsion pairs, \\ and relative periodicity. Version arXiv:2509.07645v4 (CC BY 4.0), distributed by arXiv under the Creative Commons Attribution 4.0 International licence, http://creativecommons.org/licenses/by/4.0/, as recorded in the arXiv licence metadata for this exact version and shown on the abstract page https://arxiv.org/abs/2509.07645v4. Full licence notice: \texttt{licenses/arxiv-2509.07645-CC-BY-4.0.txt}.\par\smallskip

\noindent\textit{Editorial scope.} Counted unit: the article's lem ac=ac-bctr-equivalent-conditions (source0.tex lines 3027--3042). The article's complete original proof of it is delivered with it (source0.tex lines 3044--3113, one \texttt{proof} environment of 70 lines). No mathematics is written, repaired or re-derived: the statement and the proof are the article's own lines, in the article's own notation, and no retained line is substituted or re-flowed.  Retained uncounted as the source-local context the proof needs: lines 934--942; lines 1343--1354; lines 1894--1899; lines 2491--2499; lines 2871--2884.  Nothing was omitted from any retained span, so the package carries no omission record.  No retained line contains a citation, so the wrapper carries no bibliography and no bibliography entry.  No retained line contains a figure environment, an image-inclusion command or drawing code, so no image is delivered and the package declares no asset dependency.  Editorial formatting only, in addition: an AMS \texttt{amsart} preamble that restates the article's own declarations and macros used by the retained lines, namely the environments \texttt{lem}, \texttt{prop}, \texttt{cor} on separate counters, together with the standard packages those lines need.  The retained environments print in this document as Lemma 1, Proposition 1, Corollary 1 in order of appearance, whereas the article prints the same items with the numbering of its own full document.  The complete native source of the article as distributed in the arXiv e-print for this exact version is bundled unchanged as \texttt{sources/184756/source0.tex}.
\begin{lem} \label{Ext1-as-homotopy-Hom-lemma}
 Let\/ $\sK$ be an exact category, and let $A^\bu$ and $B^\bu$ be two
complexes in\/~$\sK$.
 Assume that\/ $\Ext^1_\sK(A^n,B^n)=0$ for all $n\in\boZ$.
 Then there is a natural isomorphism of abelian groups
$$
 \Ext^1_{\Com(\sK)}(A^\bu,B^\bu)\simeq\Hom_{\Hot(\sK)}(A^\bu,B^\bu[1]).
$$
\end{lem}

\begin{prop} \label{all-contraacyclic-cotorsion-pair-prop}
 Let\/ $\sK$ be a Grothendieck category and $(\sA,\sB)$ be
a hereditary complete cotorsion pair generated by a set of objects\/
$\sS$ in\/~$\sK$.
 Then the two classes\/ $\sA''=\Com(\sA)\subset\Com(\sK)$ and $\sB''=
\Ac^\bctr(\sB)\subset\Com(\sB)\subset\Com(\sK)$ form a hereditary
complete cotorsion pair $(\sA'',\sB'')$ generated by a set of objects
in the Grothendieck category\/ $\Com(\sK)$.
 The kernel\/ $\sA''\cap\sB''$ of the cotorsion pair $(\sA'',\sB'')$
in\/ $\Com(\sK)$ is the class of all contractible complexes with
the terms in\/ $\sA\cap\sB$.
\end{prop}

\begin{cor} \label{contraacyclic-in-B-are-acyclic-in-B}
 Let\/ $\sK$ be a Grothendieck category and $(\sA,\sB)$ be a hereditary
complete cotorsion pair generated by a set of objects in\/~$\sK$.
 Then every Becker-contraacyclic complex in\/ $\sB$ is acyclic in\/
$\sB$, that is\/ $\Ac^\bctr(\sB)\subset\Ac(\sB)$.
\end{cor}

\begin{lem} \label{finite-coresol-dim-periodicity}
 Let\/ $\sK$ be an idempotent-complete exact category and\/
$\CA\subset\sK$ be a coresolving full subcategory.
 Assume that there is a finite integer~$d$ such that the coresolution
dimensions of all the objects of\/ $\sK$ with respect to\/ $\CA$ do
not exceed~$d$.
 Then every complex in\/ $\CA$ that is acyclic in\/ $\sK$ is also
acyclic in\/~$\CA$.
\end{lem}

\begin{lem} \label{proj-resolving-inj-coresolving}
\textup{(a)} Let\/ $\sE$ be an exact category and\/ $\sA\subset\sE$
be a resolving full subcategory closed under direct summands in\/~$\sE$.
 Then the classes of projective objects in the exact categories\/ $\sE$
and\/ $\sA$ coincide, $\sA_\proj=\sE_\proj$.
 There are enough projective objects in\/ $\sA$ if and only if
there are enough projective objects in\/~$\sE$. \par
\textup{(b)} Let\/ $\sE$ be an exact category and\/ $\sB\subset\sE$ be
a coresolving full subcategory closed under direct summands in\/~$\sE$.
 Then the classes of injective objects in the exact categories\/ $\sE$
and\/ $\sB$ coincide, $\sB^\inj=\sE^\inj$.
 There are enough injective objects in\/ $\sB$ if and only if
there are enough injective objects in\/~$\sE$.
\end{lem}

\begin{lem} \label{ac=ac-bctr-equivalent-conditions}
 Let\/ $\sK$ be a Grothendieck category, $(\VF,\CA)$ be a cotorsion pair
of the very flat type in\/ $\sK$, and\/ $\Flat\subset\sK$ be a resolving
full subcategory closed under direct summands in\/ $\sK$ such that\/
$\VF\subset\Flat$.
 Then the following two conditions are equivalent:
\begin{enumerate}
\item for any given complex $F^\bu$ in\/ $\Flat\cap\CA$ such that
$F^\bu$ is acyclic in\/ $\sK$, the objects of cocycles of $F^\bu$ belong
to\/ $\Flat$ if and only if, for every complex $P^\bu$ in\/ $\VF$,
every morphism of complexes $P^\bu\rarrow F^\bu$ is homotopic to zero;
\item the classes of acyclic and Becker-contraacyclic complexes in
the exact category\/ $\Flat\cap\CA$ coincide, i.~e.,
$\Ac^\bctr(\Flat\cap\CA)=\Ac(\Flat\cap\CA)$.
\end{enumerate}
\end{lem}

\begin{proof}
 There are three observations underlying the assertion of the lemma:
\begin{enumerate}
\renewcommand{\theenumi}{\Alph{enumi}}
\item A complex $F^\bu$ in $\Flat\cap\CA$ is acyclic in
$\Flat\cap\CA$ if and only if $F^\bu$ is acyclic in $\sK$ with
the objects of cocycles belonging to~$\Flat$.
 The point is that any complex in $\CA$ that is acyclic in $\sK$ is
also acyclic in $\CA$ (i.~e., has the objects of cocycles belonging
to~$\CA$), by Lemma~\ref{finite-coresol-dim-periodicity}.
\item A complex $F^\bu$ in $\Flat\cap\CA$ is Becker-contraacyclic
in $\Flat\cap\CA$ if and only if, for every complex $P^\bu$ in $\VF$,
every morphism of complexes $P^\bu\rarrow F^\bu$ is homotopic to zero.
 Indeed, by Proposition~\ref{all-contraacyclic-cotorsion-pair-prop} and
Lemma~\ref{Ext1-as-homotopy-Hom-lemma}, a complex $K^\bu$ in $\CA$ is
Becker-contraacyclic in $\CA$ if and only if, for every complex $P^\bu$
in $\VF$, every morphism of complexes $P^\bu\rarrow K^\bu$ is homotopic
to zero.
 It remains to point out that $\Flat\cap\CA$ is a resolving full
subcategory closed under direct summands in $\CA$, hence the classes of
projective objects in $\CA$ and $\Flat\cap\CA$ coincide by
Lemma~\ref{proj-resolving-inj-coresolving}(a), and a complex in
$\Flat\cap\CA$ is Becker-contraacyclic in $\Flat\cap\CA$ if and only
if it is Becker-contraacyclic in~$\CA$.
\item Every Becker-contraacyclic complex in $\Flat\cap\CA$ is
acyclic in~$\sK$.
 Indeed, every Becker-contraacyclic complex in $\CA$ is acyclic in
$\CA$ by Corollary~\ref{contraacyclic-in-B-are-acyclic-in-B}.
\end{enumerate}

 Having collected the observations above, let us spell out the proof
of the promised equivalence (1)~$\Longleftrightarrow$~(2).

 (1)~$\Longrightarrow$~(2) Let $F^\bu$ be an acyclic complex in
the exact category $\Flat\cap\CA$.
 Then, in particular, $F^\bu$ is an acyclic complex in $\sK$ with
the objects of cocycles belonging to $\Flat$.
 By~(1), it follows that, for every complex $P^\bu$ in $\VF$,
every morphism of complexes $P^\bu\rarrow F^\bu$ is homotopic to zero.
 Using~(B), we see that the complex $F^\bu$ is Becker-contraacyclic
in $\Flat\cap\CA$.

 Conversely, let $F^\bu$ be a Becker-contraacyclic complex in
the exact category $\Flat\cap\CA$.
 By~(B), it follows that, for any complex $P^\bu$ in $\VF$, every
morphism of complexes $P^\bu\rarrow F^\bu$ is homotopic to zero.
 By~(C), the complex $F^\bu$ is acyclic in~$\sK$.
 By~(1), it follows that the objects of cocycles of $F^\bu$ belong
to $\Flat$.
 Using~(A), we can conclude that the complex $F^\bu$ is acyclic
in $\Flat\cap\CA$.

 (2)~$\Longrightarrow$~(1) Let $F^\bu$ be a complex in $\Flat\cap\CA$
such that $F^\bu$ is acyclic in $\sK$ with the objects of cocycles
belonging to $\Flat$.
 By~(A), it follows that the complex $F^\bu$ is acyclic in
$\Flat\cap\CA$.
 By~(2), the complex $F^\bu$ is Becker-contraacyclic in $\Flat\cap\CA$.
 Using~(B), we see that, for every complex $P^\bu$ in $\VF$, every
morphism of complexes $P^\bu\rarrow F^\bu$ is homotopic to zero.

 Conversely, let $F^\bu$ be a complex in $\Flat\cap\CA$ such that
$F^\bu$ is acyclic in $\sK$ and, for every complex $P^\bu$ in $\VF$,
every morphism of complexes $P^\bu\rarrow F^\bu$ is homotopic to zero.
 By~(B), it follows that the complex $F^\bu$ is Becker-contraacyclic
in $\Flat\cap\CA$.
 By~(2), the complex $F^\bu$ is acyclic in $\Flat\cap\CA$.
 In particular, it follows that the objects of cocycles of $F^\bu$
belong to $\Flat$.
\end{proof}

\end{document}
